\section{Proofs and Supporting Results}\label{sec:proofs_ch4}

This section provides standard learning-theoretic helper results and the proofs for the results presented in this chapter.

\subsection{Supporting Results}


In this section, we establish the supporting results required for the proofs of the theorems and lemmas presented in the chapter.

We now state the theorem characterizing generalization bounds via Rademacher complexity.

\begin{theorem}[Rademacher Complexity Generalization Bound, \citealp{shalevshwartz2014understanding}]\label{thm:rademacher_gen}
Let $\mathcal{F}$ be a family of real-valued functions mapping from a domain $\mathcal{Z}$ to a bounded interval $[a, b]$, and let $S = (z_1, \dots, z_n) \sim \mathcal{D}^n$ be an i.i.d. sample of size $n$. Then, for any $\delta \in (0,1)$, with probability at least $1-\delta$ over the random draw of $S$, we have
\begin{equation*}
    \sup_{f \in \mathcal{F}} \left| \E[f] - \widehat{\E}_S[f] \right| \le 2\mathrm{Rad}_n(\mathcal{F}) + (b-a)\sqrt{\frac{\log(2/\delta)}{2n}},
\end{equation*}
where $\widehat{\E}_S[f] := \frac{1}{n}\sum_{i=1}^n f(z_i)$ is the empirical average over $S$, $\widehat{\mathrm{Rad}}_S(\mathcal{F}) := \E_{\sigma}\left[ \sup_{f \in \mathcal{F}} \frac{1}{n}\sum_{i=1}^n \sigma_i f(z_i) \right]$ is the empirical Rademacher complexity on sample $S$, and $\mathrm{Rad}_n(\mathcal{F}) := \E_S[\widehat{\mathrm{Rad}}_S(\mathcal{F})]$ is the expected Rademacher complexity over sample size $n$.
\end{theorem}
\begin{proof}
Let $g(S) := \sup_{f \in \mathcal{F}} (\E[f] - \widehat{\E}_S[f])$. Changing any single sample point $z_i \to z_i'$ alters $g(S)$ by at most $(b-a)/n$. Applying McDiarmid's concentration inequality (Lemma 26.4 of \citealp{shalevshwartz2014understanding}) yields $g(S) \le \E_S[g(S)] + (b-a)\sqrt{\frac{\log(2/\delta)}{2n}}$ with probability at least $1-\delta/2$. By Lemma 26.2 of \citealp{shalevshwartz2014understanding}, the expected supremum satisfies $\E_S[g(S)] \le 2 \E_S[\widehat{\mathrm{Rad}}_S(\mathcal{F})] = 2\mathrm{Rad}_n(\mathcal{F})$. Applying the identical concentration bound to the negated class $-\mathcal{F} := \{-f : f \in \mathcal{F}\}$ (noting $\mathrm{Rad}_n(-\mathcal{F}) = \mathrm{Rad}_n(\mathcal{F})$) bounds the lower-tail supremum $\sup_{f \in \mathcal{F}} (\widehat{\E}_S[f] - \E[f]) = \sup_{h \in -\mathcal{F}} (\E[h] - \widehat{\E}_S[h])$ by the same threshold with confidence $1-\delta/2$. By a union bound over both tail events holding with probability at least $1-\delta$, the identity $\sup_{f \in \mathcal{F}} |\E[f] - \widehat{\E}_S[f]| = \max\left\{ \sup_{f \in \mathcal{F}} (\E[f] - \widehat{\E}_S[f]), \, \sup_{f \in \mathcal{F}} (\widehat{\E}_S[f] - \E[f]) \right\}$ yields the two-sided absolute supremum bound.
\end{proof}

We now establish a useful property regarding the Rademacher complexity of a product function class, expressed in terms of the individual function classes.
\begin{lemma}[Rademacher Complexity of Product Function Classes]\label{lem:product_rademacher}
Let $\mathcal{F}_1$ and $\mathcal{F}_2$ be two function classes mapping from domain $\mathcal{Z}$ to $\mathbb{R}$, bounded almost surely by $M_1 = \sup_{f_1 \in \mathcal{F}_1} \|f_1\|_\infty$ and $M_2 = \sup_{f_2 \in \mathcal{F}_2} \|f_2\|_\infty$ respectively. Define the product function class $\mathcal{F}_1 \cdot \mathcal{F}_2 := \{ z \mapsto f_1(z) f_2(z) : f_1 \in \mathcal{F}_1, f_2 \in \mathcal{F}_2 \}$. Then, for any sample $S = (z_1, \dots, z_n)$, the empirical Rademacher complexity satisfies
\begin{equation*}
    \widehat{\mathrm{Rad}}_S(\mathcal{F}_1 \cdot \mathcal{F}_2) \le M_2 \widehat{\mathrm{Rad}}_S(\mathcal{F}_1) + M_1 \widehat{\mathrm{Rad}}_S(\mathcal{F}_2),
\end{equation*}
and taking expectation over $S \sim \mathcal{D}^n$ yields $\mathrm{Rad}_n(\mathcal{F}_1 \cdot \mathcal{F}_2) \le M_2 \mathrm{Rad}_n(\mathcal{F}_1) + M_1 \mathrm{Rad}_n(\mathcal{F}_2)$.
\end{lemma}
\begin{proof}
Fix a sample $S = (z_1, \dots, z_n) \in \mathcal{Z}^n$ and pick an arbitrary reference function $h_1 \in \mathcal{F}_1$. For any $f_1 \in \mathcal{F}_1$ and $f_2 \in \mathcal{F}_2$, writing $f_1(z_i) f_2(z_i) = (f_1(z_i) - h_1(z_i)) f_2(z_i) + h_1(z_i) f_2(z_i)$ and applying sub-additivity of the supremum yields
\begin{align*}
    \widehat{\mathrm{Rad}}_S(\mathcal{F}_1 \cdot \mathcal{F}_2) &\le \E_\sigma \left[ \sup_{f_1 \in \mathcal{F}_1, f_2 \in \mathcal{F}_2} \frac{1}{n}\sum_{i=1}^n \sigma_i (f_1(z_i) - h_1(z_i)) f_2(z_i) \right] \\
    &\qquad + \E_\sigma \left[ \sup_{f_2 \in \mathcal{F}_2} \frac{1}{n}\sum_{i=1}^n \sigma_i h_1(z_i) f_2(z_i) \right].
\end{align*}
For the first term, since $|f_2(z_i)| \le M_2$, the scalar mapping $t \mapsto t \cdot f_2(z_i)$ is $M_2$-Lipschitz. Applying Talagrand's contraction lemma (Lemma 26.9 of \citealp{shalevshwartz2014understanding}) bounds this term by
\begin{equation*}
    M_2 \E_\sigma \left[ \sup_{f_1 \in \mathcal{F}_1} \frac{1}{n}\sum_{i=1}^n \sigma_i (f_1(z_i) - h_1(z_i)) \right] = M_2 \widehat{\mathrm{Rad}}_S(\mathcal{F}_1),
\end{equation*}
where the linear term $\frac{1}{n}\sum_{i=1}^n \sigma_i h_1(z_i)$ vanishes under expectation since $\E_{\sigma_i}[\sigma_i] = 0$.

For the second term, since $|h_1(z_i)| \le M_1$, the scalar mapping $t \mapsto h_1(z_i) \cdot t$ is $M_1$-Lipschitz. Applying Talagrand's contraction lemma bounds this term by
\begin{equation*}
    M_1 \E_\sigma \left[ \sup_{f_2 \in \mathcal{F}_2} \frac{1}{n}\sum_{i=1}^n \sigma_i f_2(z_i) \right] = M_1 \widehat{\mathrm{Rad}}_S(\mathcal{F}_2).
\end{equation*}
Summing both bounds yields $\widehat{\mathrm{Rad}}_S(\mathcal{F}_1 \cdot \mathcal{F}_2) \le M_2 \widehat{\mathrm{Rad}}_S(\mathcal{F}_1) + M_1 \widehat{\mathrm{Rad}}_S(\mathcal{F}_2)$. Taking expectation over sample $S \sim \mathcal{D}^n$ establishes the population bound $\mathrm{Rad}_n(\mathcal{F}_1 \cdot \mathcal{F}_2) \le M_2 \mathrm{Rad}_n(\mathcal{F}_1) + M_1 \mathrm{Rad}_n(\mathcal{F}_2)$.
\end{proof}

We now restate the regret bound theorem for Kernelized Thompson Sampling on Reproducing Kernel Hilbert Spaces (RKHS) under sub-Gaussian noise feedback from \citet{chowdhury2017kernelized}.

\begin{theorem}[Kernelized Thompson Sampling Regret Bound, {\citealp[Theorem 4]{chowdhury2017kernelized}}]\label{thm:kernel_ts_master}
Let $\mathcal{D} \subset [0, r]^d$ be a compact and convex domain, and let the objective function $f: \mathcal{D} \to \mathbb{R}$ belong to a Reproducing Kernel Hilbert Space $\mathcal{H}_k(\mathcal{D})$ associated with kernel $k$, satisfying $\|f\|_{\mathcal{H}_k} \le B_k < \infty$. Suppose the observation noise $\eta_t = Z_t - f(a_t)$ is conditionally zero-mean and $R$-sub-Gaussian. Let $L = \sup_{x \in \mathcal{D}} \max_{j \in [d]} \left(\left. \frac{\partial^2 k(p,q)}{\partial p_j \partial q_j} \right|_{p=q=x}\right)^{1/2} < \infty$. Running GP-TS with round-$t$ decision sets $\mathcal{D}_t \subset \mathcal{D}$ satisfying $|f(x) - f([x]_t)| \le 1/t^2$ for all $x \in \mathcal{D}$ (via a discretization of mesh size $\epsilon_t \le \frac{1}{B_k L t^2}$) guarantees that, with probability at least $1 - \delta$, the cumulative regret satisfies
\begin{equation*}
    \operatorname{Reg}_T = \cO\left( \sqrt{(\gamma_T + \log(2/\delta)) d \log(B_k d T)} \left(\sqrt{T \gamma_T} + B_k \sqrt{T \log(2/\delta)}\right) \right) = \widetilde{\cO}\left( d^{1/2} \gamma_T \sqrt{T} \right),
\end{equation*}
where $\gamma_T$ is the maximum information gain after $T$ rounds for kernel $k$ on domain $\mathcal{D}$.
\end{theorem}
\begin{proof}
This is Theorem 4 of \citet{chowdhury2017kernelized}. We refer the reader to \citet{chowdhury2017kernelized} for the full martingale concentration derivation.
\end{proof}



\subsection{Proofs for the Decision Problem}

In this subsection, we present the proofs of the theoretical results for the decision problem.

We now present the proof of Lemma~\ref{lemma:regression_vanilla_risk}.

\begin{proof}(Proof of Lemma \ref{lemma:regression_vanilla_risk})\phantomsection\label{proof:vanilla_risk}\label{app:proof_vanilla_risk}
Under mandatory disclosure, the sender always reveals features ($s(U,X) = x$). The features $X$ are successfully received with probability $1 - q$ and are erased with probability $q$. If erased, the receiver predicts the default value $0$ and suffers quadratic loss $Y^2$. If received, the receiver predicts $b^\top X$ and suffers $(Y - b^\top X)^2$. The vanilla risk is
\begin{equation*}
    \RV(b) = q\E[Y^2] + (1 - q)\E[(Y - b^\top X)^2].
\end{equation*}
Since $\E[Y] = 0$, we have $\E[Y^2] = \mathrm{Var}(Y)$.
To decompose the second term, we add and subtract the OLS predictor $b_V^\top X$:
\begin{equation*}
    \E[(Y - b^\top X)^2] = \E\left[ \left( (Y - b_V^\top X) + (b_V - b)^\top X \right)^2 \right].
\end{equation*}
Expanding the square yields
\begin{equation*}
    \E[(Y - b^\top X)^2] = \E[(Y - b_V^\top X)^2] + \E[((b_V - b)^\top X)^2] + 2\E[(Y - b_V^\top X)(b_V - b)^\top X].
\end{equation*}
By the orthogonality principle of linear least squares, the error $Y - b_V^\top X$ is orthogonal to the feature space $\mathcal{X}$, meaning $\E[(Y - b_V^\top X)X] = 0$. Consequently, the cross-product term is zero:
\begin{equation*}
    2(b_V - b)^\top \E[(Y - b_V^\top X)X] = 0.
\end{equation*}
Furthermore, expanding the OLS residual risk yields
\begin{equation*}
    \E[(Y - b_V^\top X)^2] = \E[Y^2] - 2\E[Y b_V^\top X] + \E[(b_V^\top X)^2].
\end{equation*}
Applying orthogonality to $b_V^\top X$ gives $\E[(Y - b_V^\top X) b_V^\top X] = b_V^\top \E[(Y - b_V^\top X) X] = 0$, which implies $\E[Y b_V^\top X] = \E[(b_V^\top X)^2] = \mathrm{Var}(b_V^\top X)$. Since $\E[Y] = 0$ and $\E[b_V^\top X] = 0$, we have
\begin{equation*}
    \E[(Y - b_V^\top X)^2] = \mathrm{Var}(Y) - 2\mathrm{Var}(b_V^\top X) + \mathrm{Var}(b_V^\top X) = \mathrm{Var}(Y) - \mathrm{Var}(b_V^\top X),
\end{equation*}
and
\begin{equation*}
    \E[((b_V - b)^\top X)^2] = \mathrm{Var}((b_V - b)^\top X) = \mathrm{Var}((b - b_V)^\top X).
\end{equation*}
Substituting these back into the risk formula gives
\begin{align*}
    \RV(b) &= q\mathrm{Var}(Y) + (1 - q)\left( \mathrm{Var}(Y) - \mathrm{Var}(b_V^\top X) + \mathrm{Var}((b - b_V)^\top X) \right) \nonumber \\
    &= \mathrm{Var}(Y) - (1 - q)\mathrm{Var}(b_V^\top X) + (1 - q)\mathrm{Var}((b - b_V)^\top X),
\end{align*}
which matches~\eqref{eq:vanilla_risk_decomp}. 

The vanilla risk is minimized when $\mathrm{Var}((b - b_V)^\top X) = 0$. Assuming $\mathrm{Var}(X) \succ 0$, this occurs uniquely at $b = b_V$, yielding the optimal vanilla risk
\begin{equation*}
    \RV^*(q) = \mathrm{Var}(Y) - (1 - q)\mathrm{Var}(b_V^\top X).
\end{equation*}
\end{proof}

We now present the proof of Lemma~\ref{lem:best_response}.

\begin{proof}(Proof of Lemma \ref{lem:best_response})\phantomsection\label{proof:best_response}\label{app:proof_best_response}
Consider a sender with preference realization $u$ and features $x$ under a committed receiver predictor parameter $b$. 
If the sender chooses to withhold ($s(u,x) = \emptyset$), the receiver predicts $0$, and the sender suffers loss:
\begin{equation*}
    \ell_S(0, u) = u^2.
\end{equation*}
If the sender chooses to disclose ($s(u,x) = x$), the feature is successfully transmitted with probability $1 - q$ (resulting in prediction $b^\top x$ and loss $(b^\top x - u)^2$) and is erased with probability $q$ (resulting in prediction $0$ and loss $u^2$). The sender's expected loss from disclosing is
\begin{equation*}
    \E_{\epsilon}[\ell_S(\widehat{Y}, u)] = (1 - q)(b^\top x - u)^2 + q u^2.
\end{equation*}
The sender prefers disclosure if and only if:
\begin{equation*}
    u^2 > (1 - q)(b^\top x - u)^2 + q u^2.
\end{equation*}
Since $q \in (0, 1)$, we divide by the positive constant $1 - q > 0$ to get
\begin{equation*}
    u^2 > (u - b^\top x)^2.
\end{equation*}
Expanding the right-hand side gives
\begin{equation*}
    u^2 > u^2 - 2u b^\top x + (b^\top x)^2 \iff (2u - b^\top x)b^\top x > 0.
\end{equation*}
Recalling the definition of $F_b(u, x)$, this simplifies to $F_b(u, x) > 0$. We resolve ties by withholding, which completes the proof.
\end{proof}

We now present the proof of Lemma~\ref{lemma:regression_strategic_risk}.

\begin{proof}(Proof of Lemma \ref{lemma:regression_strategic_risk})\phantomsection\label{proof:strategic_risk}\label{app:proof_strategic_risk}
If $s_b(U,X) = x$, the sender attempts revelation. With probability $q$, the channel erases the features, leading to prediction $0$ and loss $Y^2$. With probability $1 - q$, features are successfully received, resulting in prediction $b^\top X$ and loss $(Y - b^\top X)^2$. If $s_b(U,X) = \emptyset$, the feature is withheld, leading to prediction $0$ and loss $Y^2$.
The strategic risk is
\begin{equation*}
    \RS(b) = \E\left[ \I{s_b(U,X) = \emptyset} Y^2 + \I{s_b(U,X) = x} \left( q Y^2 + (1 - q)(Y - b^\top X)^2 \right) \right].
\end{equation*}
Grouping the terms containing $Y^2$ yields
\begin{align*}
    \RS(b) &= \E\left[ Y^2 - (1 - q)\I{s_b(U,X) = x} Y^2 + (1 - q)\I{s_b(U,X) = x} (Y - b^\top X)^2 \right] \nonumber \\
    &= \E\left[ Y^2 \right] - (1 - q) \E\left[ \I{s_b(U,X) = x} \left( Y^2 - (Y - b^\top X)^2 \right) \right].
\end{align*}
Since $\E[Y] = 0$, we have $\E[Y^2] = \mathrm{Var}(Y)$. Expanding the term $Y^2 - (Y - b^\top X)^2$ gives:
\begin{equation*}
    Y^2 - (Y^2 - 2Y b^\top X + (b^\top X)^2) = (2Y - b^\top X)b^\top X = G_b(X, Y).
\end{equation*}
Thus, the expression reduces to:
\begin{equation*}
    \RS(b) = \mathrm{Var}(Y) - (1 - q)\E\left[ G_b(X, Y) \I{F_b(U, X) > 0} \right],
\end{equation*}
which is the stated equation~\eqref{eq:strategic_risk_formula}.
\end{proof}

We now present the proof of Theorem~\ref{thm:sufficient_condition}.

\begin{proof}(Proof of Theorem \ref{thm:sufficient_condition})\phantomsection\label{proof:sufficient_condition}\label{app:proof_sufficient_condition}
Let $\Delta(b) = \RV^* - \RS(b)$. By definition, if there exists $b \in \R^d$ such that $\Delta(b) > 0$, then $\Delta(\env) > 0$. Using the risk equations from Lemmas~\ref{lemma:regression_vanilla_risk} and~\ref{lemma:regression_strategic_risk}, the risk difference is
\begin{equation*}
    \RS(b) - \RV^* = (1 - q) \left[ \mathrm{Var}(b_V^\top X) - \E\left[ G_b(X, Y) \I{F_b(U, X) > 0} \right] \right].
\end{equation*}
We rewrite the expectation as
\begin{equation*}
    \E\left[ G_b(X, Y) \I{F_b(U, X) > 0} \right] = \E[G_b(X, Y)] - \E\left[ G_b(X, Y) \I{F_b(U, X) \le 0} \right].
\end{equation*}
Note that $G_b(X, Y) = (2Y - b^\top X)b^\top X = Y^2 - (Y - b^\top X)^2$. Taking expectations and using the vanilla risk calculations from Lemma~\ref{lemma:regression_vanilla_risk}, we have
\begin{align*}
    \E[G_b(X, Y)] &= \E[Y^2] - \E[(Y - b^\top X)^2] \nonumber \\
    &= \mathrm{Var}(Y) - \left( \mathrm{Var}(Y) - \mathrm{Var}(b_V^\top X) + \mathrm{Var}((b - b_V)^\top X) \right) \nonumber \\
    &= \mathrm{Var}(b_V^\top X) - \mathrm{Var}((b - b_V)^\top X).
\end{align*}
Thus, the risk difference is
\begin{equation*}
    \RS(b) - \RV^* = (1 - q) \left[ \mathrm{Var}((b - b_V)^\top X) + \E\left[ G_b(X, Y) \I{F_b(U, X) \le 0} \right] \right].
\end{equation*}
For this difference to be strictly negative, we require the bracketed term to be negative.
We decompose the expectation term
\begin{equation*}
    \E\left[ G_b \I{F_b \le 0} \right] = \E\left[ G_b \I{G_b \le 0} \right] + \E\left[ G_b \left( \I{F_b \le 0} - \I{G_b \le 0} \right) \right].
\end{equation*}
For the first term, we use the standard algebraic identity $z \I{z \le 0} = \frac{z - |z|}{2}$ for any real scalar $z$. Applying this to $G_b(X, Y)$ and substituting $\E[G_b(X, Y)] = \mathrm{Var}(b_V^\top X) - \mathrm{Var}((b - b_V)^\top X)$ yields
\begin{align}
    \E\left[ G_b \I{G_b \le 0} \right] &= \frac{\E[G_b] - \E[|G_b|]}{2} \nonumber \\
    &= \frac{\mathrm{Var}(b_V^\top X) - \mathrm{Var}((b - b_V)^\top X) - \E[|G_b(X,Y)|]}{2}.
\end{align}
For the second term, we establish the pointwise upper bound 
\begin{equation*}
    G_b \left( \I{F_b \le 0} - \I{G_b \le 0} \right) \le 2|Y - U||b^\top X|.
\end{equation*}
Hence:
\begin{equation*}
    G_b \left( \I{F_b \le 0} - \I{G_b \le 0} \right) \le |G_b - F_b| = 2|Y - U||b^\top X|.
\end{equation*}
Taking expectation and applying the Cauchy-Schwarz inequality gives:
\begin{equation*}
    \E\left[ G_b \left( \I{F_b \le 0} - \I{G_b \le 0} \right) \right] \le 2\E\left[ |Y - U||b^\top X| \right] \le 2 \|U - Y\|_2 \sqrt{\mathrm{Var}(b^\top X)}.
\end{equation*}
Combining these components, we get
\begin{equation*}
    \E\left[ G_b \I{F_b \le 0} \right] \le \frac{\mathrm{Var}(b_V^\top X) - \mathrm{Var}((b - b_V)^\top X) - \E[|G_b|]}{2} + 2 \|U - Y\|_2 \sqrt{\mathrm{Var}(b^\top X)}.
\end{equation*}
Substitute this inequality back into the risk difference:
\begin{equation*}
    \RS(b) - \RV^* \le (1 - q) \left[ \frac{\mathrm{Var}(b_V^\top X) + \mathrm{Var}((b - b_V)^\top X) - \E[|G_b|]}{2} + 2\|U - Y\|_2 \sqrt{\mathrm{Var}(b^\top X)} \right].
\end{equation*}
To guarantee that $\RS(b) - \RV^* < 0$, it is sufficient for the term in the brackets to be strictly negative:
\begin{equation*}
    2\|U - Y\|_2 \sqrt{\mathrm{Var}(b^\top X)} < \frac{\E[|G_b|] - \mathrm{Var}(b_V^\top X) - \mathrm{Var}((b - b_V)^\top X)}{2}.
\end{equation*}
Rearranging this inequality yields the sufficient condition.
\end{proof}

We now present the proof of Theorem~\ref{thm:necessary_condition}.

\begin{proof}(Proof of Theorem \ref{thm:necessary_condition})\phantomsection\label{proof:necessary_condition}\label{app:proof_necessary_condition}
We analyze the risk difference:
\begin{equation*}
    \RS(b) - \RV^* = (1 - q) \left[ \mathrm{Var}((b - b_V)^\top X) + \E\left[ G_b(X, Y) \I{F_b(U, X) \le 0} \right] \right].
\end{equation*}
Pointwise, for any realizations, we have
\begin{equation*}
    G_b \I{F_b \le 0} \ge G_b \I{G_b \le 0}.
\end{equation*}

Taking expectations on both sides yields
\begin{equation*}
    \E\left[ G_b \I{F_b \le 0} \right] \ge \E\left[ G_b \I{G_b \le 0} \right] = \frac{\mathrm{Var}(b_V^\top X) - \mathrm{Var}((b - b_V)^\top X) - \E[|G_b|]}{2}.
\end{equation*}
Substitute this lower bound into the risk difference:
\begin{align*}
    \RS(b) - \RV^* &\ge (1 - q) \left[ \mathrm{Var}((b - b_V)^\top X) + \frac{\mathrm{Var}(b_V^\top X) - \mathrm{Var}((b - b_V)^\top X) - \E[|G_b|]}{2} \right] \nonumber \\
    &= \frac{1 - q}{2} \left[ \mathrm{Var}(b_V^\top X) + \mathrm{Var}((b - b_V)^\top X) - \E[|G_b|] \right].
\end{align*}
Now, suppose that for all $b \in \R^d$, the necessary condition fails, meaning:
\begin{equation*}
    \E[|G_b(X, Y)|] - \mathrm{Var}(b_V^\top X) - \mathrm{Var}((b - b_V)^\top X) \le 0 \quad \forall b \in \R^d.
\end{equation*}
This implies that the term inside the brackets is always non-negative:
\begin{equation*}
    \mathrm{Var}(b_V^\top X) + \mathrm{Var}((b - b_V)^\top X) - \E[|G_b(X, Y)|] \ge 0 \quad \forall b \in \R^d.
\end{equation*}
Since $1 - q > 0$, this guarantees $\RS(b) - \RV^* \ge 0$ for all $b$, which means $\Delta(\env) \le 0$.
\end{proof}


\subsection{Proofs for Offline Learning}

In this subsection, we provide the proofs of the generalization bounds for offline learning.

We now present the proof of Lemma~\ref{lem:vanilla_generalization}.

\begin{proof}(Proof of Lemma \ref{lem:vanilla_generalization})\phantomsection\label{proof:vanilla_generalization}\label{app:proof_vanilla_generalization}
Let $L(b) = \E[(Y - b^\top X)^2]$ denote the expected quadratic loss under receiver linear strategy $b \in \mathcal{R}$, and let $\widehat{L}(b) = \frac{1}{n}\sum_{k=1}^n (y_k - b^\top x_k)^2$ denote the empirical loss on sample $S$. The difference between true and empirical vanilla risks is
\begin{align*}
    \RV(b) - \RhatV(b) &= q\left( \mathrm{Var}(Y) - \frac{1}{n}\sum_{k=1}^n y_k^2 \right) + (1 - q)(L(b) - \widehat{L}(b)).
\end{align*}
Applying the triangle inequality yields
\begin{equation*}
    \sup_{b \in \mathcal{R}} |\RV(b) - \RhatV(b)| \le \left| \mathrm{Var}(Y) - \frac{1}{n}\sum_{k=1}^n y_k^2 \right| + \sup_{b \in \mathcal{R}} |L(b) - \widehat{L}(b)|.
\end{equation*}
We allocate failure probability $\delta/2$ to each of these two terms using a union bound.

For the first term, since $Y \in [-1, 1]$ almost surely, $Y^2 \in [0, 1]$. Applying Hoeffding's concentration inequality yields $\left| \mathrm{Var}(Y) - \frac{1}{n}\sum_{k=1}^n y_k^2 \right| \le \sqrt{\frac{\log(4/\delta)}{2n}} = \cO\left( \sqrt{\frac{\log(1/\delta)}{n}} \right)$ with probability at least $1 - \delta/2$.

To bound the second term $\sup_{b \in \mathcal{R}} |L(b) - \widehat{L}(b)|$, we apply the Rademacher complexity generalization bound (Theorem~\ref{thm:rademacher_gen}) to the loss-composed strategy space $\ell \circ \mathcal{R} := \{ (x,y) \mapsto (y - b^\top x)^2 : b \in \mathcal{R} \}$. First, we show that the scalar loss function is Lipschitz over the domain of predictions. For any strategy $b \in \mathcal{R}$ (where $\|b\|_\infty \le 1$) and feature vector $x$ (where $\|x\|_\infty \le 1$), the absolute prediction is bounded by $|b^\top x| = |\sum_{i=1}^d b_i x_i| \le \sum_{i=1}^d |b_i||x_i| \le d \cdot 1 = d$, so the scalar linear prediction $a = b^\top x$ lies in $[-d, d]$. Define the scalar loss function $\phi(a) = (y - a)^2$ for $a \in [-d, d]$ and $y \in [-1, 1]$. Differentiating $\phi(a)$ with respect to $a$ gives $|\phi'(a)| = 2|y - a| \le 2(1 + d)$. Hence, $\phi$ is $L$-Lipschitz over $[-d, d]$ with Lipschitz constant $L = 2(1 + d)$.

Next, by Talagrand's contraction lemma (Lemma 26.9 of \citealp{shalevshwartz2014understanding}), the empirical Rademacher complexity of the loss-composed class $\ell \circ \mathcal{R}$ is bounded by
\begin{equation*}
    \widehat{\mathrm{Rad}}_S(\ell \circ \mathcal{R}) \le L \cdot \widehat{\mathrm{Rad}}_S(\mathcal{R}) = 2(1+d) \widehat{\mathrm{Rad}}_S(\mathcal{R}),
\end{equation*}
where $\widehat{\mathrm{Rad}}_S(\mathcal{R})$ is the empirical Rademacher complexity of the receiver strategy space $\mathcal{R} = \{ x \mapsto b^\top x : \|b\|_\infty \le 1 \}$. By Lemma 26.10 of \citet{shalevshwartz2014understanding} (Rademacher complexity bound for linear classes), for any sample $S = (x_1, \dots, x_n)$ with $\|x_k\|_2 \le \sqrt{d}$ and $\|b\|_2 \le \sqrt{d}$:
\begin{equation*}
    \widehat{\mathrm{Rad}}_S(\mathcal{R}) = \E_\sigma \left[ \sup_{b \in \mathcal{R}} \frac{1}{n}\sum_{k=1}^n \sigma_k b^\top x_k \right] \le \frac{\|b\|_2 \max_k \|x_k\|_2}{\sqrt{n}} \le \frac{\sqrt{d} \cdot \sqrt{d}}{\sqrt{n}} = \frac{d}{\sqrt{n}}.
\end{equation*}
Thus, $\widehat{\mathrm{Rad}}_S(\ell \circ \mathcal{R}) \le 2(1+d) \frac{d}{\sqrt{n}} = \frac{2d(1+d)}{\sqrt{n}} = \cO\left( \frac{d^2}{\sqrt{n}} \right)$, and taking expectation over $S \sim \mathcal{D}^n$ gives $\mathrm{Rad}_n(\ell \circ \mathcal{R}) = \cO\left( \frac{d^2}{\sqrt{n}} \right)$.

Finally, applying Theorem~\ref{thm:rademacher_gen} to class $\ell \circ \mathcal{R}$ (bounded in $[0, (1+d)^2]$), with probability at least $1 - \delta/2$:
\begin{equation*}
    \sup_{b \in \mathcal{R}} |L(b) - \widehat{L}(b)| \le 2\mathrm{Rad}_n(\ell \circ \mathcal{R}) + (1+d)^2 \sqrt{\frac{\log(4/\delta)}{2n}} = \cO\left( \frac{d^2}{\sqrt{n}} + d^2\sqrt{\frac{\log(1/\delta)}{n}} \right).
\end{equation*}
Summing the bounds for both terms via a union bound completes the proof.
\end{proof}

We now present the proof of Lemma~\ref{lem:strategic_generalization}.

\begin{proof}(Proof of Lemma \ref{lem:strategic_generalization})\phantomsection\label{proof:strategic_generalization}\label{app:proof_strategic_generalization}
\textit{Notation abuse: Except for $R_S$ and $\widehat{R}_S$, where the subscript $_S$ denotes \emph{strategic}, the subscript $_S$ elsewhere denotes the underlying sample set $S$.}

\noindent The strategic risk and empirical strategic risk are
\begin{equation*}
    \RS(b) = \mathrm{Var}(Y) - (1 - q)\E[h_b(U, X, Y)], \;\; \widehat{\RS}(b) = \frac{1}{n}\sum_{k=1}^n y_k^2 - (1 - q)\frac{1}{n}\sum_{k=1}^n h_b(u_k, x_k, y_k),
\end{equation*}
where $h_b(u, x, y) = G_b(x, y)\I{F_b(u, x) > 0}$. Applying the triangle inequality:
\begin{equation*}
    \sup_{b \in \mathcal{R}} |\RS(b) - \widehat{\RS}(b)| \le \left| \mathrm{Var}(Y) - \frac{1}{n}\sum_{k=1}^n y_k^2 \right| + (1-q) \sup_{b \in \mathcal{R}} \left|\E[h_b] - \widehat{\E}_S[h_b]\right|.
\end{equation*}
The first term is bounded by $\cO\left(\sqrt{\log(1/\delta)/n}\right)$ with probability at least $1 - \delta/2$ using Hoeffding's inequality.

To bound the second term $\sup_{b \in \mathcal{R}} |\E[h_b] - \widehat{\E}_S[h_b]|$, consider the function class $\mathcal{H} = \{ h_b(u, x, y) = G_b(x, y)\I{F_b(u, x) > 0} : b \in \mathcal{R} \}$. Note that each $h_b = g_b \cdot k_b$ is a pointwise product of quadratic polynomial $g_b(x, y) = (2y - b^\top x)b^\top x$ from class $\mathcal{G}_1 = \{ (x,y) \mapsto (2y - b^\top x)b^\top x : b \in \mathcal{R} \}$ and binary threshold indicator $k_b(u, x) = \I{(2u - b^\top x)b^\top x > 0}$ from class $\mathcal{K} = \{ (u,x) \mapsto \I{(2u - b^\top x)b^\top x > 0} : b \in \mathcal{R} \}$. Since $\|b\|_\infty \le 1$, $\|x\|_\infty \le 1$, and $|y| \le 1$, by Hölder's inequality $|b^\top x| \le d$, which implies $\|g_b\|_\infty \le (2 + d)d = d(d+2)$ and $\|k_b\|_\infty \le 1$. By Lemma~\ref{lem:product_rademacher}, the empirical Rademacher complexity of $\mathcal{H} = \mathcal{G}_1 \cdot \mathcal{K}$ satisfies
\begin{equation*}
    \begin{aligned}
    \widehat{\mathrm{Rad}}_S(\mathcal{H}) &\le \sup_{b \in \mathcal{R}} \|g_b\|_\infty \cdot \widehat{\mathrm{Rad}}_S(\mathcal{K}) + \sup_{b \in \mathcal{R}} \|k_b\|_\infty \cdot \widehat{\mathrm{Rad}}_S(\mathcal{G}_1) \\
    &\le d(d+2) \widehat{\mathrm{Rad}}_S(\mathcal{K}) + \widehat{\mathrm{Rad}}_S(\mathcal{G}_1).
    \end{aligned}
\end{equation*}

For the quadratic term $\mathcal{G}_1$, write $g_b(x_k, y_k) = \phi_k(b^\top x_k)$ where $\phi_k(t) = 2y_k t - t^2$ is a $1$-variable $L$-Lipschitz scalar function with $L = 2(1+d)$ over $t \in [-d, d]$. Applying Talagrand's contraction lemma (Lemma 26.9 of \citealp{shalevshwartz2014understanding}) to $\widehat{\mathrm{Rad}}_S(\mathcal{G}_1) = \E_\sigma \left[ \sup_{b \in \mathcal{R}} \frac{1}{n}\sum_{k=1}^n \sigma_k \phi_k(b^\top x_k) \right]$ factors out the Lipschitz constant $L$
\begin{equation*}
    \widehat{\mathrm{Rad}}_S(\mathcal{G}_1) \le 2(1+d) \widehat{\mathrm{Rad}}_S(\mathcal{R}) \le 2(1+d)\frac{d}{\sqrt{n}} = \cO\left( \frac{d^2}{\sqrt{n}} \right),
\end{equation*}
where $\widehat{\mathrm{Rad}}_S(\mathcal{R}) \le \frac{d}{\sqrt{n}}$ is the linear predictor bound (Lemma 26.10 of \citealp{shalevshwartz2014understanding}).

For the indicator class $\mathcal{K}$, each function $\I{F_b(u, x) > 0}$ is a thresholded quadratic polynomial $F_b(u, x) = 2 u (b^\top x) - (b^\top x)^2$ in variables $(u, x) \in \R^{d+1}$. The vector space $V$ of degree-2 polynomials in $d+1$ variables has dimension $\dim(V) = \frac{(d+3)(d+2)}{2} = \cO(d^2)$. By Theorem 3.5 of \citet{anthony1999neural} (VC-dimension of thresholded vector spaces of functions), $\mathrm{VC}(\mathcal{K}) \le \dim(V) = \frac{(d+3)(d+2)}{2} = \cO(d^2)$. For sample $S = (z_1, \dots, z_n)$, the finite set of binary evaluation vectors $A_S = \{ (k(z_1), \dots, k(z_n)) \in \{0, 1\}^n : k \in \mathcal{K} \}$ has size bounded by the growth function $\tau_{\mathcal{K}}(n) := \max_{S'} |A_{S'}|$. By the Sauer-Shelah-Perles Lemma (Lemma 6.10 of \citealp{shalevshwartz2014understanding}), $|A_S| \le \tau_{\mathcal{K}}(n) \le (en / \mathrm{VC}(\mathcal{K}))^{\mathrm{VC}(\mathcal{K})}$. Since each binary vector $a \in A_S$ has $L_2$-norm $\|a\|_2 = \sqrt{\sum_{i=1}^n a_i^2} \le \sqrt{n}$, applying Massart's Finite Class Lemma (Lemma 26.8 of \citealp{shalevshwartz2014understanding}) to $A_S$ yields
\begin{equation*}
    \begin{aligned}
    \widehat{\mathrm{Rad}}_S(\mathcal{K}) &= \widehat{\mathrm{Rad}}_S(A_S) = \E_\sigma \left[ \frac{1}{n}\sup_{a \in A_S} \sum_{i=1}^n \sigma_i a_i \right] \le \frac{\sqrt{n}\sqrt{2 \log |A_S|}}{n} \\
    &\le \sqrt{\frac{2 \mathrm{VC}(\mathcal{K}) \log(en / \mathrm{VC}(\mathcal{K}))}{n}} \le \sqrt{\frac{2 \mathrm{VC}(\mathcal{K}) \log(en)}{n}} = \cO\left( \frac{d \sqrt{\log n}}{\sqrt{n}} \right).
    \end{aligned}
\end{equation*}

Combining both bounds gives $\widehat{\mathrm{Rad}}_S(\mathcal{H}) \le d(d+2) \cO\left( \frac{d \sqrt{\log n}}{\sqrt{n}} \right) + \cO\left( \frac{d^2}{\sqrt{n}} \right) = \cO\left( \frac{d^3 \sqrt{\log n}}{\sqrt{n}} \right)$, and taking expectation yields $\mathrm{Rad}_n(\mathcal{H}) = \cO\left( \frac{d^3 \sqrt{\log n}}{\sqrt{n}} \right)$.

Finally, applying Theorem~\ref{thm:rademacher_gen} to class $\mathcal{H}$ (bounded by $\|h_b\|_\infty \le d(d+2)$), with probability at least $1 - \delta/2$
\begin{equation*}
    \begin{aligned}
    \sup_{b \in \mathcal{R}} |\E[h_b] - \widehat{\E}_S[h_b]| &\le 2\mathrm{Rad}_n(\mathcal{H}) + d(d+2)\sqrt{\frac{\log(4/\delta)}{2n}} \\
    &= \cO\left( \frac{d^3 \sqrt{\log n}}{\sqrt{n}} + d^2 \sqrt{\frac{\log(1/\delta)}{n}} \right).
    \end{aligned}
\end{equation*}
Summing these bounds via a union bound completes the proof.
\end{proof}

We now present the proof of Theorem~\ref{thm:advantage_generalization}.

\begin{proof}(Proof of Theorem \ref{thm:advantage_generalization})\phantomsection\label{proof:advantage_generalization}\label{app:proof_advantage_generalization}
Let $\Delta(\env) = \min_{b \in \mathcal{R}}\RV(b) - \min_{b \in \mathcal{R}}\RS(b)$ and $\widehat{\Delta}(D_n) = \min_{b \in \mathcal{R}} \widehat{\RV}(b) - \min_{b \in \mathcal{R}} \widehat{\RS}(b)$. We bound the difference using
\begin{equation*}
    \left| \min_{b \in \mathcal{R}} f(b) - \min_{b \in \mathcal{R}} g(b) \right| \le \sup_{b \in \mathcal{R}} |f(b) - g(b)|.
\end{equation*}
\begin{align*}
    |\Delta(\env) - \widehat{\Delta}(D_n)| &\le \left| \min_{b \in \mathcal{R}} \RV(b) - \min_{b \in \mathcal{R}} \widehat{\RV}(b) \right| + \left| \min_{b \in \mathcal{R}} \RS(b) - \min_{b \in \mathcal{R}} \widehat{\RS}(b) \right| \nonumber \\
    &\le \sup_{b \in \mathcal{R}} |\RV(b) - \widehat{\RV}(b)| + \sup_{b \in \mathcal{R}} |\RS(b) - \widehat{\RS}(b)|.
\end{align*}
Using a union bound, both generalization events hold simultaneously with probability at least $1 - \delta$. Summing the bounds yields the advantage generalization bound.
\end{proof}


\subsection{Proofs for Online Learning}

In this subsection, we provide the proofs for the online learning problem.

We now present the proof of Lemma~\ref{lem:sobolev_membership}.

\begin{proof}(Proof of Lemma \ref{lem:sobolev_membership})\phantomsection\label{proof:sobolev_membership}\label{app:proof_sobolev_membership}
We establish Sobolev space membership by showing that both the expected vanilla risk $\RV(b)$ and strategic risk $\RS(b)$ are infinitely differentiable ($\mathcal{C}^\infty$) on the compact domain $\mathcal{R} \subset \mathbb{R}^d$, which would imply $L(m,b)$ is infinite differentiable on compact domain which further implies their Sobolev norms of any order are finite.

\noindent\textbf{Vanilla Risk $\RV(b)$:}
Recall the formula for the vanilla risk
\begin{equation*}
    \RV(b) = \mathrm{Var}(Y) - (1-q)\mathrm{Var}(b_V^\top X) + (1-q)\mathrm{Var}((b-b_V)^\top X).
\end{equation*}
Since $\RV(b)$ is a quadratic polynomial in $b$, all its partial derivatives of order $|\alpha| \ge 3$ vanish identically. Thus, $\RV(b) \in \mathcal{C}^\infty(\mathcal{R})$.

\noindent\textbf{Strategic Risk $\RS(b)$ Smoothness on $\mathcal{R}$:}
Recall the strategic risk formula:
\begin{equation*}
    \RS(b) = \mathrm{Var}(Y) - (1-q)\E[G_b(X, Y)\I{F_b(U, X) > 0}],
\end{equation*}
where $F_b(u, x) = (2u - b^\top x) b^\top x$ and $G_b(x, y) = 2y b^\top x - (b^\top x)^2$. We show that $I(b) := \E[G_b(X, Y)\I{F_b(U, X) > 0}] \in \mathcal{C}^\infty(\mathcal{R})$.

For any $b \in \mathcal{R}$ with $\|b\|_2 \ge b_{\min} > 0$, perform the elementary change of variables $v = u - \frac{1}{2}b^\top x$. The revelation condition $(2u - b^\top x)b^\top x > 0$ simplifies to $v(b^\top x) > 0$. Let $\hat{b} = b / \|b\|_2$, and let $R_b \in \mathrm{SO}(d)$ be a smooth local rotation mapping $\hat{b}$ to the first standard basis vector $e_1$. In rotated coordinates $z = R_b x$ (so $x = R_b^\top z$), the first coordinate satisfies $z_1 = \hat{b}^\top x$, which implies $b^\top x = \|b\|_2 z_1$.

Because $\|b\|_2 > 0$, the revelation condition $v(b^\top x) > 0$ is equivalent to $v z_1 > 0$, which defines the fixed region
\begin{equation*}
    \Omega_0 := \{(v, z_1) \in \mathbb{R}^2 : v z_1 > 0\} \times \mathbb{R}^{d-1} \times \mathbb{R}.
\end{equation*}
Crucially, the domain of integration $\Omega_0$ has zero dependence on $b$. Since the Jacobian of the transformation $(u, x, y) \mapsto (v, z, y)$ is identically $1$, change of variables yields
\begin{equation*}
    I(b) = \int_{\Omega_0} G_b(R_b^\top z, y) \, p_{U,X,Y}\left(v + \tfrac{1}{2}\|b\|_2 z_1,\, R_b^\top z,\, y\right) dv \, dz \, dy.
\end{equation*}
Because $p_{U,X,Y}$ is $\mathcal{C}^\infty$ with compact support, $G_b$ is a polynomial in $(z, y, b)$, and the mapping $b \mapsto (\|b\|_2, R_b)$ is smooth for all $\|b\|_2 \ge b_{\min} > 0$, the integrand is an infinitely differentiable function of $b$ whose partial derivatives of all orders are uniformly bounded on the fixed integration domain $\Omega_0$. Applying the standard Leibniz integral rule (dominated convergence theorem) directly justifies differentiating under the integral sign to any order $k$. Therefore, $I(b) \in \mathcal{C}^\infty(\mathcal{R})$, and consequently $\RS(b) \in \mathcal{C}^\infty(\mathcal{R})$.

% =========================================================================
% Original Inductive Differentiation Proof (Preserved for Reference)
% =========================================================================
\iffalse
We show by induction on $k = |\alpha|$ that for any multi-index $\alpha \in \mathbb{N}_0^d$ of order $k$, the partial derivative $D^\alpha I(b)$ exists, is continuous on $\mathcal{R}$, and can be expressed as a finite linear combination of terms belonging to two canonical integral forms
\begin{equation}
\label{eq:form_a}
\begin{aligned}
    \text{Form A: } \quad \int_{\mathbb{R}^d \times \mathbb{R}} P(x,y,b) \Big[ &\I{b^\top x > 0}\int_{\frac{1}{2}b^\top x}^{\infty} D_u^{\alpha_3} p \, du \\
    + &\I{b^\top x < 0}\int_{-\infty}^{\frac{1}{2}b^\top x} D_u^{\alpha_3} p \, du \Big] dy \, dx,
\end{aligned}
\end{equation}
\begin{equation}
\label{eq:form_b}
    \text{Form B: } \quad \int_{\mathbb{R}^d \times \mathbb{R}} Q(x,y,b) D_x^{\alpha_1} D_y^{\alpha_2} D_u^{\alpha_3} p\left(x, y, \tfrac{1}{2}b^\top x\right) dy \, dx,
\end{equation}
where $P(x,y,b)$ and $Q(x,y,b)$ are polynomials in $(x,y,b)$, and $D_x^{\alpha_1} D_y^{\alpha_2} D_u^{\alpha_3} p$ represents partial derivatives of $p$ of order $|\alpha_1|+|\alpha_2|+\alpha_3 \le k$.

\noindent\textbf{Base Case ($k=1$):} Differentiating $I(b)$ with respect to $b_i$ via Leibniz's rule yields
\begin{align*}
    \frac{\partial I(b)}{\partial b_i} &= \int_{\mathbb{R}^d \times \mathbb{R}} \frac{\partial G_b(x, y)}{\partial b_i} \left[ \I{b^\top x > 0} \int_{\frac{1}{2}b^\top x}^{\infty} p \, du + \I{b^\top x < 0} \int_{-\infty}^{\frac{1}{2}b^\top x} p \, du \right] dy \, dx \\
    &\quad + \int_{\mathbb{R}^d \times \mathbb{R}} \frac{1}{2}x_i G_b(x, y) \text{sign}(-b^\top x) p\left(x, y, \tfrac{1}{2}b^\top x\right) dy \, dx.
\end{align*}
The first integral is of Form A (with $P = \frac{\partial G_b}{\partial b_i}$ and $\alpha_3=0$), and the second integral is of Form B (with $Q = \frac{1}{2}x_i G_b$ and $\alpha_3=0$). Note that the discontinuity introduced by the indicator function does not contribute, since $G_b$ vanishes on the boundary ${x : b^\top x = 0}$.

\noindent\textbf{Inductive Step ($k \to k+1$):} Assume that $D^\alpha I(b)$ is a finite sum of terms of Form A and Form B for all $|\alpha| \le k$. Differentiating a generic term of Form A with respect to $b_j$ via Leibniz's rule yields
\begin{align*}
    \frac{\partial}{\partial b_j} (\text{Form A}) &= \int_{\mathbb{R}^d \times \mathbb{R}} \frac{\partial P(x,y,b)}{\partial b_j} \Big[ \I{b^\top x > 0} \int_{\frac{1}{2}b^\top x}^{\infty} D_u^{\alpha_3} p \, du \\
    &+ \I{b^\top x < 0} \int_{-\infty}^{\frac{1}{2}b^\top x} D_u^{\alpha_3} p \, du \Big] dy \, dx \\
    &\quad + \int_{\mathbb{R}^d \times \mathbb{R}} P(x,y,b) \Big[ \text{sign}(-b^\top x) \frac{1}{2}x_j D_u^{\alpha_3} p\left(x,y,\tfrac{1}{2}b^\top x\right)\Big] dy \, dx.
\end{align*}

This expansion directly decomposes into two structurally identical forms:
(a) A term of Form A where the polynomial factor becomes $\tilde{P}(x,y,b) = \frac{\partial P(x,y,b)}{\partial b_j}$.
(b) A boundary term of Form B with the associated polynomial factor $\tilde{Q}(x,y,b) = \text{sign}(-b^\top x) \frac{1}{2}x_j P(x,y,b)$ determined by the active indicator region.

Similarly, differentiating a generic term of Form B with respect to $b_j$ via the product and chain rules yields
\begin{align*}
    \frac{\partial}{\partial b_j} (\text{Form B}) &= \int_{\mathbb{R}^d \times \mathbb{R}} \frac{\partial Q(x,y,b)}{\partial b_j} D_x^{\alpha_1} D_y^{\alpha_2} D_u^{\alpha_3} p\left(x, y, \tfrac{1}{2}b^\top x\right) dy \, dx \\
    &\quad + \int_{\mathbb{R}^d \times \mathbb{R}} Q(x,y,b) \left( \frac{1}{2}x_j \right) D_x^{\alpha_1} D_y^{\alpha_2} D_u^{\alpha_3+1} p\left(x, y, \tfrac{1}{2}b^\top x\right) dy \, dx.
\end{align*}
This calculation again produces a combination of terms that fit our canonical structure:
(a) A term of Form B where the polynomial factor updates to $\tilde{Q}_1(x,y,b) = \frac{\partial Q(x,y,b)}{\partial b_j}$, preserving the derivative order of the density.
(b) Another term of Form B where the new polynomial factor is $\tilde{Q}_2(x,y,b) = \frac{1}{2}x_j Q(x,y,b)$, and the order of the density derivative with respect to $u$ increases to $\alpha_3+1$. 

Thus it follows that $\RS(b) \in \mathcal{C}^\infty(\mathcal{R})$.
\fi

% \noindent\textbf{Alternate Proof of Smoothness via Distributional Derivatives:}
% The smoothness of the strategic risk can also be established more elegantly using the theory of distributions. The indicator functions $\mathbb{I}\{b^\top x > 0\}$ and $\mathbb{I}\{b^\top x < 0\}$ can be viewed as Heaviside step functions parameterized by $b$. In the sense of distributions, differentiating these step functions with respect to $b$ yields Dirac delta functions supported on the boundary hyperplane $b^\top x = 0$. Higher-order derivatives yield derivatives of the Dirac delta. 

% The integral $I(b)$ can be interpreted as the action of these parameterized distributions on the test functions involving $p_{U,X,Y}(u,x,y)$ and $G_b(x,y)$. By the fundamental property of distributional derivatives, evaluating the derivative of a distribution against a test function is equivalent to evaluating the distribution against the derivative of the test function (i.e., integration by parts). Because the density $p_{U,X,Y}$ is infinitely differentiable ($\mathcal{C}^\infty$) and has compact support, it (along with the polynomial $G_b$) forms a valid smooth test function that can safely absorb any number of transferred derivatives. Therefore, applying the derivative to the integral always results in a bounded, well-defined functional evaluation. This immediately guarantees that the integral $I(b)$, and consequently $\RS(b)$, is infinitely differentiable with respect to $b$, bypassing the need for explicit boundary term tracking.

\noindent\textbf{Sobolev Space Membership:}
Since $\mathcal{R} \subset \mathbb{R}^d$ is compact and both $\RV(b), \RS(b) \in \mathcal{C}^\infty(\mathcal{R})$, every partial derivative $D^\alpha L(m,b)$ is continuous on the compact set $\mathcal{R}$ and therefore bounded: $\sup_{b \in \mathcal{R}} |D^\alpha L(m,b)| \le M_\alpha < \infty$.
Consequently, the Sobolev norm of order $m_0 = 5/2 + d/2$ is finite:
\begin{align*}
    \|L(m, \cdot)\|_{H^{5/2+d/2}(\mathcal{R})}^2 &= \sum_{|\alpha| \le \lceil 5/2 + d/2 \rceil} \int_{\mathcal{R}} |D^\alpha L(m, b)|^2 \, db \\
    &\le \operatorname{Vol}(\mathcal{R}) \sum_{|\alpha| \le \lceil 5/2 + d/2 \rceil} M_\alpha^2 < \infty.
\end{align*}
Thus, $L(m, b) \in H^{5/2+d/2}(\mathcal{R})$.
\end{proof}

We now present the proof of Theorem~\ref{thm:kernel_online_regret}.

\begin{proof}(Proof of Theorem \ref{thm:kernel_online_regret})\phantomsection\label{proof:kernel_online_regret}\label{app:proof_kernel_online_regret}
The proof proceeds by verifying the premises required to apply the kernelized Thompson Sampling regret bound of \citet[Theorem 4]{chowdhury2017kernelized} (stated as Theorem~\ref{thm:kernel_ts_master}) to our strategic setting over $\mathcal{A} = \{V, S\} \times \mathcal{R}$.

\noindent\textbf{1. Noise Sub-Gaussianity:} 
Because the domain of features, outcomes, and decision vectors is bounded ($\|X\|_2 \le X_0, |Y| \le Y_0, \|b\|_2 \le 1$), the realized loss $Z_t = \ell(m_t, b_t, X_t, Y_t, U_t)$ is uniformly bounded in a finite interval $[0, M_0]$. For any joint action $a_t = (m_t, b_t)$, both $Z_t$ and its conditional expectation $\mathbb{E}[Z_t \mid a_t]$ are bounded in $[0, M_0]$. By Hoeffding's Lemma, the zero-mean observation noise $\eta_t = Z_t - \mathbb{E}[Z_t \mid a_t]$ is bounded in $[-M_0, M_0]$ and thus conditionally $R$-sub-Gaussian with $R = M_0/2$.

\noindent\textbf{2. RKHS Membership and Direct-Sum Norm Bound:} 
By Lemma~\ref{lem:sobolev_membership}, the expected risk functions under each regime belong to the Sobolev space $H^{5/2+d/2}(\mathcal{R})$, which is norm-equivalent to the RKHS $\mathcal{H}_{k_0}(\mathcal{R})$ of the Matérn-$5/2$ kernel: $\|R_V\|_{\mathcal{H}_{k_0}} \le B_{k,V} < \infty$ and $\|R_S\|_{\mathcal{H}_{k_0}} \le B_{k,S} < \infty$. For the direct-sum kernel $k((m, b), (m', b')) = \mathbb{I}\{m = m'\} k_0(b, b')$, the joint RKHS decomposes orthogonally as $\mathcal{H}_k = \mathcal{H}_{k_0} \oplus \mathcal{H}_{k_0}$. Consequently, the joint objective $f(m, b) = \mathbb{I}\{m=V\} R_V(b) + \mathbb{I}\{m=S\} R_S(b)$ belongs to $\mathcal{H}_k(\mathcal{A})$ with bounded joint RKHS norm
\begin{equation*}
    \|f\|_{\mathcal{H}_k} = \sqrt{\|R_V\|_{\mathcal{H}_{k_0}}^2 + \|R_S\|_{\mathcal{H}_{k_0}}^2} \le \sqrt{B_{k,V}^2 + B_{k,S}^2} \le \sqrt{2}\max\{B_{k,V}, B_{k,S}\} =: B_k < \infty.
\end{equation*}

\noindent\textbf{3. Lipschitz Regularity and Action Discretization Net:} 
For the Matérn-$5/2$ kernel, $k_0$ is twice continuously differentiable on the compact set $\mathcal{R}$, so $L := \sup_{b \in \mathcal{R}} \max_{j \in [d]} \left(\left. \frac{\partial^2 k_0(p,q)}{\partial p_j \partial q_j}\right|_{p=q=b}\right)^{1/2} < \infty$. By standard RKHS properties \citep[Eq.~(7)]{chowdhury2017kernelized}, any $f \in \mathcal{H}_k$ is Lipschitz with constant $L_f \le B_k L$. Evaluating candidate actions over the round-$t$ decision set $\mathcal{A}_t = \{V, S\} \times \mathcal{R}_t$, where $\mathcal{R}_t$ is an $\epsilon_t$-net of $\mathcal{R}$ with mesh size $\epsilon_t \le \frac{1}{B_k L t^2}$, guarantees that $|f(a) - f([a]_t)| \le B_k L \epsilon_t \le 1/t^2$ for all $a \in \mathcal{A}$, satisfying the exact decision-set premise of Theorem~\ref{thm:kernel_ts_master}. The cumulative discretization error across all $T$ rounds is bounded by $\sum_{t=1}^T 1/t^2 \le \pi^2/6 = \mathcal{O}(1)$.

\noindent\textbf{4. Regret Derivation via Information Gain:} 
Applying Theorem~\ref{thm:kernel_ts_master} with sub-Gaussian noise parameter $R$, RKHS norm bound $B_k$, and round-$t$ decision net $\mathcal{A}_t$ guarantees that, with probability at least $1 - \delta$, the cumulative regret satisfies
\begin{equation*}
    \operatorname{Reg}_T = \cO\left( d^{1/2} \gamma_T \sqrt{T \log(dT/\delta)} \right).
\end{equation*}
For the direct-sum Matérn-$5/2$ kernel across the two regimes, the maximum information gain satisfies $\gamma_T \le 2\gamma_{T, \mathcal{R}} = \mathcal{O}\left(T^{\frac{d}{5+d}} (\log T)^{\frac{5}{5+d}}\right)$ \citep[Table 1]{vakili2021information}. Substituting this rate yields
\begin{equation*}
    \operatorname{Reg}_T = \cO\left( d^{1/2} T^{\frac{5+3d}{10+2d}} (\log T)^{\frac{5}{5+d}} \sqrt{\log(dT/\delta)} \right) = \widetilde{\mathcal{O}}\left( T^{\frac{5+3d}{10+2d}} \right).
\end{equation*}
For any dimension $d \in \{1, 2, 3, 4\}$, the exponent $\frac{5+3d}{10+2d} < 1$, guaranteeing sublinear cumulative regret $\operatorname{Reg}_T = o(T)$.
\end{proof}

\iffalse
% Previous proof version preserved for reference:
\begin{proof}(Proof of Theorem \ref{thm:kernel_online_regret})
The proof proceeds by verifying the conditions required to apply the continuous Gaussian Process bandit regret bound of \citet{chowdhury2017kernelized} to our strategic setting.

We begin by establishing the sub-Gaussianity of the feedback noise. Because the domain of features, outcomes, and decision vectors is bounded ($\|X\|_2 \le X_0, |Y| \le Y_0, \|b\|_2 \le 1$), the realized loss $Z_t = \ell(m_t, b_t, X_t, Y_t, U_t)$ is uniformly bounded in a finite interval $[0, M_0]$. For the joint action $a_t = (m_t, b_t)$, both $Z_t$ and its conditional expectation $\mathbb{E}[Z_t \mid a_t]$ are bounded in $[0, M_0]$. By Hoeffding's Lemma, the zero-mean observation noise $\eta_t = Z_t - \mathbb{E}[Z_t \mid a_t]$ is bounded in $[-M_0, M_0]$ and thus conditionally $R$-sub-Gaussian with $R = M_0/2$.

Next, we establish RKHS membership. By the definition of $\mathcal{P}_{\text{Sob}}$ (Definition~\ref{def:sobolev_risk}), the expected loss $L(m, b)$, which evaluates to $\RV(b)$ or $\RS(b)$ depending on the regime $m$, belongs to the Sobolev space $H^{5/2+d/2}(\mathcal{R})$. It is a standard result that the Reproducing Kernel Hilbert Space (RKHS) of the Matérn-$5/2$ kernel is norm-equivalent to the Sobolev space $H^{5/2+d/2}(\mathcal{R})$. Therefore, the joint objective $f(m, b)$ successfully lies within the product Matérn-$5/2$ RKHS $\mathcal{H}_k$ over $\mathcal{A}$ with a bounded RKHS norm $\|f\|_{\mathcal{H}_k} \le B_k := \max\{B_{k,V}, B_{k,S}\} < \infty$.

Finally, we combine these properties to derive the regret bound. With $R$-sub-Gaussian noise and RKHS membership $f \in \mathcal{H}_k$ with $\|f\|_{\mathcal{H}_k} \le B_k$, applying Theorem 4 of \citet{chowdhury2017kernelized} (stated as Theorem~\ref{thm:kernel_ts_master}) guarantees that GP-TS achieves cumulative regret:
\begin{equation*}
    \operatorname{Reg}_T = \cO\left( d^{1/2} \gamma_T \sqrt{T \log(dT/\delta)} \right).
\end{equation*}
For the Matérn-$5/2$ kernel (smoothness parameter $\nu = 5/2$) on a $d$-dimensional compact space, the maximum information gain is bounded by $\gamma_T = \mathcal{O}\left(T^{\frac{d}{5+d}} (\log T)^{\frac{5}{5+d}}\right)$ \citep[Table 1]{vakili2021information}. Substituting this into the regret expression $\operatorname{Reg}_T = \cO\left( d^{1/2} \gamma_T \sqrt{T \log(dT/\delta)} \right)$ yields the final explicit bound:
\begin{equation*}
    \operatorname{Reg}_T = \cO\left( d^{1/2} T^{\frac{5+3d}{10+2d}} (\log T)^{\frac{5}{5+d}} \sqrt{\log(dT/\delta)} \right) = \widetilde{\mathcal{O}}\left( T^{\frac{5+3d}{10+2d}} \right).
\end{equation*}
\end{proof}
\fi

% Applied Fix: Replaced product-class Rademacher contraction bound with direct vector contraction.
